
\chapter{Zahlen -- Von 1 bis 10}

\begin{itemize}[leftmargin=*]
\item \arword{واحد}{wáaHid}{1}{\%D9\%88\%D8\%A7\%D8\%AD\%D8\%AF}
\item \arword{اثنين}{itnéin}{2}{\%D8\%A7\%D8\%AB\%D9\%86\%D9\%8A\%D9\%86}
\item \arword{ثلاثة}{taláata}{3}{\%D8\%AB\%D9\%84\%D8\%A7\%D8\%AB\%D8\%A9}
\item \arword{أربعة}{árba3a}{4}{\%D8\%A3\%D8\%B1\%D8\%A8\%D8\%B9\%D8\%A9}
\item \arword{خمسة}{khámsa}{5}{\%D8\%AE\%D9\%85\%D8\%B3\%D8\%A9}
\item \arword{ستة}{sítta}{6}{\%D8\%B3\%D8\%AA\%D8\%A9}
\item \arword{سبعة}{sáb3a}{7}{\%D8\%B3\%D8\%A8\%D8\%B9\%D8\%A9}
\item \arword{ثمانية}{tamánya}{8}{\%D8\%AB\%D9\%85\%D8\%A7\%D9\%86\%D9\%8A\%D8\%A9}
\item \arword{تسعة}{tís3a}{9}{\%D8\%AA\%D8\%B3\%D8\%B9\%D8\%A9}
\item \arword{عشرة}{3áshara}{10}{\%D8\%B9\%D8\%B4\%D8\%B1\%D8\%A9}
\end{itemize}

\begin{lachbox}
Kleine Ironie des Schicksals: die Ziffern, die wir im Deutschen
,,arabische Ziffern`` nennen (0,1,2,3...), sehen im echten Arabisch ganz
anders aus. Die Ägypter benutzen ihre eigenen Symbole -- entdecke sie unten.
\end{lachbox}

\section*{Die echten arabischen Ziffern}

\begin{center}
\renewcommand{\arraystretch}{1.8}
\begin{tabular}{cccccccccc}
\textarabic{\Large ٠} & \textarabic{\Large ١} & \textarabic{\Large ٢} &
\textarabic{\Large ٣} & \textarabic{\Large ٤} & \textarabic{\Large ٥} &
\textarabic{\Large ٦} & \textarabic{\Large ٧} & \textarabic{\Large ٨} &
\textarabic{\Large ٩} \\
0 & 1 & 2 & 3 & 4 & 5 & 6 & 7 & 8 & 9 \\
\end{tabular}
\end{center}

\begin{tippbox}
Diese Ziffern siehst du auf ägyptischen Preisschildern, Nummernschildern
und Restaurantrechnungen. Lohnt sich, sie zu erkennen, bevor du beim
Feilschen (Teil 5!) über den Preis verhandelst, den du gar nicht lesen
kannst.
\end{tippbox}

\begin{checkpointbox}
Zähle laut von 1 bis 10 auf Arabisch, während du an den Fingern
abzählst -- eine der ältesten und zuverlässigsten Lernmethoden der Welt.
\end{checkpointbox}
